% The notation (top) and the symbols of the paper (bottom); see STYLEGUIDE.md, "Mathematics".
% This paper never uses the components of a vector. So, unlike our default notation, a subscript
% numbers the elements of a series, e.g., the vertices of a triangle.
\begin{tblr}{colspec={lX[l]}, width=\linewidth}
    \toprule
    \textbf{Symbol}                 & \textbf{Meaning} \\
    \midrule
    \SetCell[c=2]{l}\textit{Notation} \\
    $a$, $\alpha$                   & scalar \\
    $\vect{a}$                      & column vector \\
    $\pnt{a}$                       & point \\
    $a_k$, $\vect{a}_k$, $\pnt{a}_k$ & $k$-th element of a series \\
    $\Vert \vect{a} \Vert$          & Euclidean norm \\
    $\vect{a}^\top$                 & transpose \\
    \midrule
    \SetCell[c=2]{l}\textit{Interpolation, \cref{sec:PerspectiveCorrect}} \\
    $\pnt{p}_k$                     & vertex $k \in \{0, 1, 2\}$ of a triangle in camera space \\
    $\vect{u}_k$                    & texture coordinates of $\pnt{p}_k$ [texels] \\
    $x_k, y_k, z_k, w_k$            & homogeneous coordinates of $\pnt{p}_k$ after the projection; $w_k$ is the depth of $\pnt{p}_k$ \\
    $\lambda_k$                     & screen-space barycentric coordinates of a pixel \\
    $\vect{u}_\mathrm{aff}$         & affinely interpolated texture coordinates \\
    $\vect{u}$                      & perspective-correct texture coordinates \\
    $\mu_k$                         & weights of perspective-correct interpolation \\
    \midrule
    \SetCell[c=2]{l}\textit{Edge and subdivision, \cref{sec:EdgeError,sec:Subdivision}} \\
    $t$                             & position of a pixel on an edge in screen space \\
    $s$                             & position of the point that the pixel shows on the edge \\
    $r$                             & depth ratio of an edge, $r = w_1 / w_0 \ge 1$ \\
    $\ell$                          & length of an edge in the texture [texels] \\
    $\delta(t)$                     & error along an edge as a fraction of $\ell$ \\
    $\delta_\mathrm{max}$, $t^\ast$ & maximum of $\delta$ and its position \\
    $n$                             & number of pieces per edge \\
    $E_n$                           & largest error after subdivision [texels] \\
    $\varepsilon$                   & error threshold [texels] \\
    \bottomrule
\end{tblr}
